\documentclass[12pt]{amsart}
\usepackage{amsmath,amssymb,amsthm}
\usepackage{hyperref}
\usepackage{tikz}

\newtheorem{theorem}{Theorem}[section]
\newtheorem{proposition}[theorem]{Proposition}
\newtheorem{lemma}[theorem]{Lemma}
\newtheorem{corollary}[theorem]{Corollary}
\newtheorem{definition}[theorem]{Definition}
\newtheorem{example}[theorem]{Example}
\newtheorem{remark}[theorem]{Remark}

\newcommand{\ZZ}{\mathbb{Z}}
\newcommand{\NN}{\mathbb{N}}
\newcommand{\Des}{\mathrm{Des}}

\title{How Atlas Enumerates the Weyl Group}
\author{}
\date{}

\begin{document}

\begin{abstract}
We explain the data structure and algorithm used by the Atlas of Lie Groups and
Representations software to represent and enumerate elements of a Weyl group.
The central idea, due to Fokko du~Cloux, is to factor each group element as a
product of minimal-length coset representatives for a chain of parabolic
subquotients, encoding the element as a fixed-length array of small integers.
Multiplication is implemented by a collection of finite-state transducer
automata, one per parabolic subquotient, that were precomputed at group
construction time.  The reference for the underlying mathematics is du~Cloux's
1999 paper~\cite{duCloux1999}.
\end{abstract}

\maketitle

\tableofcontents

%----------------------------------------------------------------------
\section{Weyl Groups and Coxeter Systems}
%----------------------------------------------------------------------

A \emph{Weyl group} $W$ is a finite group generated by reflections in a
Euclidean space; it is in particular a \emph{Coxeter group}.  We fix a set of
\emph{simple reflections} (generators) $S = \{s_1, \ldots, s_n\}$, where $n$
is the rank.  Every element of $W$ can be written as a word in these generators,
and the \emph{length} $\ell(w)$ of $w \in W$ is the minimum length of such a
word.

The generators satisfy the \emph{Coxeter relations}:
\[
  s_i^2 = e, \qquad (s_i s_j)^{m_{ij}} = e,
\]
where $m_{ij} \in \{2,3,4,6\}$ for a Weyl group (the value $m_{ij}=2$ means
$s_i$ and $s_j$ commute).  The Coxeter matrix $(m_{ij})$ encodes the same
information as the Cartan matrix and is determined by the Dynkin diagram.

The cardinalities of the Weyl groups of the simple Lie algebras are:
\[
  |W(A_n)| = (n+1)!, \quad
  |W(B_n)| = |W(C_n)| = 2^n n!, \quad
  |W(D_n)| = 2^{n-1} n!,
\]
\[
  |W(E_6)| = 51840, \quad
  |W(E_7)| = 2903040, \quad
  |W(E_8)| = 696729600,
\]
\[
  |W(F_4)| = 1152, \quad
  |W(G_2)| = 12.
\]
These grow rapidly; for example $|W(E_8)| \approx 7 \times 10^8$ and
$|W(B_{20})| = 2^{20} \cdot 20! \approx 2.7 \times 10^{24}$.  Storing all
elements as explicit reduced words or permutation matrices is infeasible.

%----------------------------------------------------------------------
\section{Parabolic Subgroups and the Subquotient Decomposition}
%----------------------------------------------------------------------

\subsection{Parabolic subgroups}

For a subset $J \subset S$, the \emph{parabolic subgroup} $W_J =
\langle s_i : s_i \in J \rangle$ is itself a Coxeter group (with the induced
Coxeter matrix), and its length function is the restriction of that of $W$.

We fix a chain
\[
  W_0 = \{e\} \subset W_1 \subset W_2 \subset \cdots \subset W_n = W,
\]
where $W_i = \langle s_1, \ldots, s_i \rangle$ (using the \emph{internal}
numbering of generators described in Section~\ref{sec:ordering}).  Since each
$W_i$ is parabolic, this is a chain of parabolic subgroups.  The index
\[
  N_i = [W_i : W_{i-1}] = \frac{|W_i|}{|W_{i-1}|}
\]
is the number of \emph{parabolic subquotient} cosets at level $i$; the group
order factors as
\[
  |W| = N_1 \cdot N_2 \cdots N_n.
\]

\subsection{Minimal-length coset representatives}

\begin{definition}
For a parabolic subgroup $W_J \subset W$, the set of
\emph{minimal-length left-coset representatives} of $W_J \backslash W$ is
\[
  {}^J W = \{w \in W : \ell(sw) > \ell(w) \text{ for all } s \in J\}.
\]
That is, ${}^J W$ consists of elements with no left descent in $J$.
\end{definition}

The fundamental theorem on parabolic cosets (see, e.g., Humphreys
\cite{Humphreys}) states:

\begin{theorem}\label{thm:coset-rep}
Every element $w \in W$ has a unique decomposition $w = u \cdot x$ with
$u \in W_J$ and $x \in {}^J W$; moreover $\ell(w) = \ell(u) + \ell(x)$.
\end{theorem}

For the restricted quotient $W_{i-1} \backslash W_i$, set $D_i = {}^{S_{i-1}}
W_i$ (where $S_{i-1} = \{s_1, \ldots, s_{i-1}\}$).  These are the elements of
$W_i$ with no left descent in $S_{i-1}$, i.e., the minimal-length
representatives of the left cosets $W_{i-1} \backslash W_i$.  Clearly
$|D_i| = N_i$.

\subsection{The iterated factorization}

Applying Theorem~\ref{thm:coset-rep} repeatedly gives the core structural
result:

\begin{theorem}[Parabolic factorization]\label{thm:factorization}
Every element $w \in W$ has a unique factorization
\[
  w = d_1 \cdot d_2 \cdots d_n
\]
with $d_i \in D_i$ for each $i$.  Moreover $\ell(w) = \sum_{i=1}^n \ell(d_i)$.
\end{theorem}

\begin{proof}
We argue by downward induction, extracting the rightmost factor first.
Since $w \in W = W_n$, Theorem~\ref{thm:coset-rep} applied to $W_{n-1} \subset
W_n$ gives a unique decomposition $w = u_{n-1} \cdot d_n$ with $u_{n-1} \in
W_{n-1}$ and $d_n \in D_n$.  Apply the same theorem to $u_{n-1} \in W_{n-1}$
and $W_{n-2} \subset W_{n-1}$ to get $u_{n-1} = u_{n-2} \cdot d_{n-1}$, and
continue.  At the last step, $u_1 \in W_1 = \langle s_1 \rangle$ decomposes as
$u_1 = d_1$ (since $D_1 = W_1$).  Uniqueness and the length equality follow
from Theorem~\ref{thm:coset-rep} at each step.
\end{proof}

\begin{example}[$W(A_2) = S_3$]\label{ex:A2}
Label the generators as in Bourbaki: $s_1 = (12)$, $s_2 = (23)$.  Then
$W_1 = \{e, s_1\}$ and $W_2 = S_3$.
\begin{itemize}
  \item $D_1 = W_1 = \{e, s_1\}$, so $N_1 = 2$.
  \item $D_2 = \{w \in W : \ell(s_1 w) > \ell(w)\} = \{e,\, s_2,\, s_2 s_1\}$,
    so $N_2 = 3$.
\end{itemize}
The six elements and their factorizations are:
\[
\begin{array}{c|cc}
w & d_1 & d_2 \\ \hline
e & e & e \\
s_1 & s_1 & e \\
s_2 & e & s_2 \\
s_1 s_2 & s_1 & s_2 \\
s_2 s_1 & e & s_2 s_1 \\
s_1 s_2 s_1 & s_1 & s_2 s_1
\end{array}
\]
Piece 1 (value 0 or 1, encoded as unsigned char 0 or 1) selects the
$W_0 \backslash W_1$ coset; piece 2 (value 0, 1, or 2) selects the
$W_1 \backslash W_2$ coset.
\end{example}

%----------------------------------------------------------------------
\section{Generator Ordering}\label{sec:ordering}
%----------------------------------------------------------------------

Atlas represents each $d_i$ by its index $a_i \in \{0, 1, \ldots, N_i - 1\}$
as an \texttt{unsigned char} (8-bit integer, values $0$--$254$; value $255$ is
reserved as a sentinel).  For this to work we need $N_i \leq 255$ for every
level $i$.

With the standard \emph{Bourbaki ordering} of generators, some classical types
fail this constraint.  For type $B_n$ with generators
$s_1 - s_2 - \cdots - s_{n-1} \Rightarrow s_n$ (double bond at the right):
\[
  W_i = \langle s_1, \ldots, s_i \rangle \cong W(A_i), \quad i < n,
  \qquad
  W_n = W(B_n).
\]
The final subquotient has size $N_n = |W(B_n)| / |W(A_{n-1})| = 2^n$, which
overflows a byte for $n \geq 8$.

The fix is to \emph{reverse} the generator order for types $B$, $C$, and $D$.
Define the internal generators $t_j = s_{n+1-j}$ (so $t_1 = s_n$, the
``special'' end of the diagram).  Then
\[
  W_r^{\mathrm{int}} = \langle t_1, \ldots, t_r \rangle \cong W(B_r),
  \qquad
  N_r = \frac{|W(B_r)|}{|W(B_{r-1})|} = 2r.
\]
The maximum subquotient size is $N_n = 2n$, which is $\leq 255$ for $n \leq
127$.

Similarly, for type $A_n$ (standard ordering) the subquotient at level $r$ has
size $N_r = r + 1$, so the maximum is $n+1 \leq 255$ for $n \leq 254$.

\begin{remark}
Atlas actually allows ranks up to \texttt{RANK\_MAX}, and the comment in the
source notes: ``the unsigned char type miraculously suffices for all subquotients
of all groups up to rank 128 (indeed, the biggest subquotient for $B_{128}$ is
of order 256), \emph{provided} the generators are enumerated in an appropriate
order.''  The reversal of $B$, $C$, $D$ is precisely this appropriate ordering.
For the exceptional types, the sizes are small enough that no reordering is
needed.
\end{remark}

Because of this reordering, the \texttt{WeylGroup} object carries two
permutations, \texttt{d\_in} and \texttt{d\_out}, that translate between the
\emph{external} Bourbaki numbering (used in all user-facing output) and the
\emph{internal} numbering (used inside the transducer tables).

%----------------------------------------------------------------------
\section{The \texttt{WeylElt} Data Type}
%----------------------------------------------------------------------

\begin{definition}
A \texttt{WeylElt} is a fixed-length array
\[
  (a_1, a_2, \ldots, a_n) \in \prod_{i=1}^n \{0, 1, \ldots, N_i - 1\},
\]
stored as \texttt{pieces[0]}, \ldots, \texttt{pieces[n-1]} of type
\texttt{unsigned char}.  Piece $i$ encodes which element of $D_{i+1}$ appears
in the factorization $w = d_1 d_2 \cdots d_n$.
\end{definition}

Since the pieces are fixed-size and small, a \texttt{WeylElt} occupies exactly
\texttt{RANK\_MAX} bytes (currently 16 bytes), regardless of the group.  This
is the same size as a 128-bit integer, so comparison and copying are extremely
cheap.

The identity element is the all-zeros tuple $(0, 0, \ldots, 0)$.  The longest
element $w_0 \in W$ has $a_i = N_i - 1$ for every $i$ (the maximal piece value
at each level).

%----------------------------------------------------------------------
\section{The Transducer Automaton}
%----------------------------------------------------------------------

\subsection{Overview}

The multiplication $w \mapsto w \cdot s$ (right multiplication by a simple
reflection) is computed using a system of precomputed finite-state automata
called \emph{transducers}, one per level.  The Transducer at level $r$ encodes
the right-multiplication action of all generators $s_1, \ldots, s_r$ on the
coset representatives $D_r$.

\subsection{Transitions and transductions}

The key algebraic fact is: for a minimal-length coset representative
$x \in D_r$ and a generator $s_j$ with $j \leq r$, exactly one of the
following holds:
\begin{enumerate}
  \item \textbf{Transition:} $x \cdot s_j = x'$ for some other $x' \in D_r$
    (with $x' \neq x$).
  \item \textbf{Transduction:} $x \cdot s_j = s_k \cdot x$ for some generator
    $s_k$ with $k < r$, i.e., $x$ is fixed and the generator $s_k$ must be
    ``passed'' to the lower-level transducers.
\end{enumerate}
In case (1), the level-$r$ piece changes; the multiplication stops at level $r$.
In case (2), the level-$r$ piece is unchanged; the emitted generator $s_k$ is
fed into the transducer at level $k' \leq r-1$ (the highest level in the same
Dynkin diagram component as $k$), and the process continues.

The two cases are distinguished by inspecting the transition table: a
\emph{shift} entry gives the new piece index $x'$, while a
\emph{transduction} entry gives the emitted generator $s_k$.  In practice
both are stored in a single table (indexed by current piece and incoming
generator), using the convention that entries $< |D_r|$ are shifts and entries
$\geq |D_r|$ encode transductions.

\subsection{Right multiplication algorithm}

Given $w = (a_1, \ldots, a_n)$ and an (internally numbered) generator $s$,
the algorithm to compute $w \cdot s$ proceeds as follows:
\begin{enumerate}
  \item Let $R$ be the index of the last generator in the Dynkin diagram
    component containing $s$ (field \texttt{upper[s]} in Atlas).  Set
    $s_{\mathrm{cur}} = s$ and $i = R$.
  \item Consult the Transducer at level $i$ with current piece $a_i$ and
    generator $s_{\mathrm{cur}}$.
  \begin{itemize}
    \item \textbf{Shift:} update $a_i \leftarrow a_i'$, stop.  Return $+1$ if
      $a_i' > a_i$ (length increases), $-1$ if $a_i' < a_i$ (length decreases).
    \item \textbf{Transduction:} emit $s_k$; set $s_{\mathrm{cur}} = s_k$
      and $i \leftarrow i - 1$; go to step 2.
  \end{itemize}
\end{enumerate}
Note that the loop can only descend to lower indices, so it terminates in at
most $n$ steps.  In practice the transduction chain is very short.

\begin{remark}
Generators in different Dynkin diagram components always commute, so right
multiplication by a generator $s$ from one component never affects pieces
belonging to a different component.  The field \texttt{upper[s]} correctly
restricts attention to the component of $s$.
\end{remark}

\subsection{Constructing the transducer tables}

The Transducer for level $r$ is built at group construction time by a
bootstrapping procedure, starting from the identity coset representative and
iteratively discovering new representatives by applying the final generator
$s_r$.

The procedure runs as follows.  Initialize the automaton with a single state
(the identity piece).  For the first state and for each state added later:
\begin{itemize}
  \item For all generators $t \leq r$: if the shift or transduction of the
    current state under $t$ is not yet defined, create a new state $xs$
    (extending the normal form by $s_r$), and determine its table entries
    by analyzing the dihedral group orbit $\langle s, t \rangle$ acting on
    the current state.
\end{itemize}
The dihedral group $\langle s, t \rangle$ has order $2m$ where $m = m_{st}$ is
the Coxeter matrix entry.  There are three cases depending on the length of the
orbit of the current element in the quotient:
\begin{enumerate}
  \item Full orbit of size $2m$: the shift target for $t$ is determined by
    traversing $m-1$ upward steps from the bottom of the orbit.
  \item Orbit of size $m$ with a stationary step at the top: a transduction
    occurs; the emitted generator is the same as that emitted by the bottom
    of the orbit under the first generator.
  \item Otherwise: the action of $t$ takes the current state upward;
    nothing needs to be recorded (the table entry remains undefined until
    a future state defines it from the other direction).
\end{enumerate}
The loop terminates because each iteration either defines a previously
undefined table entry or creates a new state (and the number of states is
bounded by $N_r \leq 255$).

%----------------------------------------------------------------------
\section{Enumeration via Mixed-Radix Representation}
%----------------------------------------------------------------------

\subsection{The bijection}

By Theorem~\ref{thm:factorization}, an element $w \in W$ is uniquely
determined by the tuple $(a_1, \ldots, a_n) \in \prod_{i=1}^n \{0, \ldots,
N_i - 1\}$.  This gives a bijection
\[
  W \leftrightarrow \{0, 1, \ldots, |W| - 1\}
\]
via the \emph{mixed-radix} (also called factorial-number system) formula
\[
  \phi(w) = a_1 + N_1 (a_2 + N_2 (a_3 + \cdots + N_{n-1} a_n \cdots)).
\]
Piece 1 is the \emph{least significant ``digit''} and piece $n$ is the most
significant.

This bijection is implemented in Atlas by two methods:
\begin{itemize}
  \item \texttt{to\_big\_int(w)}: computes $\phi(w)$ by evaluating the
    above expression using Horner's method, starting from piece $n$ and
    working down:
    \[
      u \leftarrow 0, \quad
      u \leftarrow u \cdot N_i + a_i \quad (i = n, n-1, \ldots, 1).
    \]
    The result is returned as a \texttt{big\_int} to handle groups whose order
    exceeds $2^{64}$.
  \item \texttt{to\_WeylElt(u)}: computes the inverse by successive Euclidean
    division:
    \[
      a_i \leftarrow u \bmod N_i, \quad u \leftarrow \lfloor u / N_i \rfloor
      \quad (i = 1, 2, \ldots, n).
    \]
\end{itemize}

\subsection{Iteration over all elements}

To iterate over all $|W|$ elements in order of their integer index, it
suffices to iterate through all tuples
$(a_1, a_2, \ldots, a_n)$ with $0 \leq a_i < N_i$.  In practice one can
maintain a \texttt{WeylElt} and increment it lexicographically: increment $a_1$
and carry to $a_2$ when $a_1 = N_1$, etc.  This gives the elements in an order
consistent with the Bruhat order restricted to the chain $W_1, \ldots, W_n$
(though not the full Bruhat order on $W$).

\begin{example}[Continuing Example~\ref{ex:A2}]
With $N_1 = 2$, $N_2 = 3$, the six tuples and their integer indices are:
\[
\begin{array}{c|cc|c}
w & a_1 & a_2 & \phi(w) \\ \hline
e & 0 & 0 & 0 \\
s_1 & 1 & 0 & 1 \\
s_2 & 0 & 1 & 2 \\
s_1 s_2 & 1 & 1 & 3 \\
s_2 s_1 & 0 & 2 & 4 \\
s_1 s_2 s_1 & 1 & 2 & 5
\end{array}
\]
\end{example}

%----------------------------------------------------------------------
\section{Data Structure Summary}
%----------------------------------------------------------------------

The complete picture is as follows.

\textbf{\texttt{WeylGroup} object.}  Constructed once from the Cartan matrix.
Contains:
\begin{itemize}
  \item \texttt{d\_transducer}: a vector of $n$ \texttt{Transducer} objects.
    Transducer $i$ stores a matrix \texttt{table} of size $N_{i+1} \times
    (i+2)$ (states $\times$ generators), whose entries encode transitions and
    transductions as described above.
  \item \texttt{d\_in}, \texttt{d\_out}: permutations translating between
    Bourbaki (external) and internal generator numbering.
  \item \texttt{upper}: for each (internal) generator $s$, the index of the
    last generator in the same Dynkin diagram component.
  \item \texttt{d\_min\_star}: for each generator $s$, the smallest-indexed
    generator that does not commute with $s$ (or $s$ itself if all others
    commute); used to speed up left multiplication.
\end{itemize}

\textbf{\texttt{WeylElt} value.}  A fixed array of \texttt{RANK\_MAX}
\texttt{unsigned char}s.  Piece $i$ (0-indexed) holds $a_{i+1} \in \{0,
\ldots, N_{i+1}-1\}$.  Unused pieces (for generators beyond the rank) are
held at zero.  Operations:
\begin{itemize}
  \item \textbf{Identity}: all pieces zero.
  \item \textbf{Right multiply by $s$}: \texttt{inner\_mult}, complexity
    $O(n)$ in the worst case but typically $O(1)$ or $O(2)$ in practice.
  \item \textbf{Length}: sum of the lengths of the individual pieces (stored
    in each \texttt{Transducer}).
  \item \textbf{Comparison}: lexicographic order on pieces; consistent with
    grouping by Bruhat length but not with Bruhat order itself.
  \item \textbf{Word recovery}: read off the pieces from right to left,
    recovering the canonical reduced expression.
\end{itemize}

%----------------------------------------------------------------------
\section{Size Constraints and Rank Limits}
%----------------------------------------------------------------------

The constraint that each piece fits in an \texttt{unsigned char} (values
$0$--$254$, with $255$ reserved) limits the achievable ranks:

\begin{center}
\begin{tabular}{l|l|l}
Type & Internal generator order & Max level-$r$ subquotient size \\ \hline
$A_n$ & Bourbaki & $r+1$ \\
$B_n$, $C_n$ & Reversed Bourbaki & $2r$ \\
$D_n$ & Reversed Bourbaki & $2r-2$ (for $r \geq 2$) \\
$E_6$, $E_7$, $E_8$ & Fixed (see Bourbaki) & $\leq 240$ \\
$F_4$ & Fixed & $\leq 24$ \\
$G_2$ & Fixed & $\leq 6$
\end{tabular}
\end{center}

In each case the bound is comfortably below $255$.  The combined constraint
for a level-$r$ transducer is actually $N_r + r \leq 256$: the $N_r$
shift values and (at most) $r$ possible transduction outputs must fit in
a single byte.  For $C_{86}$ this gives $171 + 85 = 256$, just at the
practical limit.

%----------------------------------------------------------------------
\section{Connection to the Bruhat Order and Kazhdan--Lusztig Theory}
%----------------------------------------------------------------------

While the parabolic factorization is designed for efficient computation rather
than reflection of Bruhat order, it does have representation-theoretic
significance.  The factorization $w = d_1 \cdots d_n$ is \emph{length-additive}:
$\ell(w) = \sum_i \ell(d_i)$.  This means that a reduced word for $w$ is
obtained by concatenating reduced words for $d_1, \ldots, d_n$ in order.

The parabolic subquotients also arise naturally when one stratifies the flag
variety $G/B$ by the Schubert cells corresponding to the chain
$B \subset P_1 \subset P_2 \subset \cdots \subset G$ (where $P_i$ is the
parabolic subgroup whose Levi factor has Weyl group $W_i$).

In Atlas, the \texttt{WeylElt} type is also used to represent
\emph{twisted involutions} (elements $\sigma \in W$ satisfying $\sigma = \delta(\sigma^{-1})$ for a diagram involution $\delta$), which parametrize the $K$-orbits on the flag variety needed for the computation of Kazhdan--Lusztig--Vogan polynomials.

%----------------------------------------------------------------------

\begin{thebibliography}{9}

\bibitem{duCloux1999}
Fokko du~Cloux,
\textit{A transducer approach to Coxeter groups},
Journal of Symbolic Computation \textbf{27} (1999), 311--324.

\bibitem{Humphreys}
James E.\ Humphreys,
\textit{Reflection Groups and Coxeter Groups},
Cambridge Studies in Advanced Mathematics, vol.\ 29,
Cambridge University Press, 1990.

\bibitem{BB}
Anders Bj\"orner and Francesco Brenti,
\textit{Combinatorics of Coxeter Groups},
Graduate Texts in Mathematics, vol.\ 231,
Springer, 2005.

\bibitem{AtlasBook}
Jeffrey Adams, Marc van Leeuwen, Peter Trapa, and David A.\ Vogan Jr.,
\textit{Unitary representations of real reductive groups},
preprint, \texttt{arXiv:1212.2192}.

\end{thebibliography}

\end{document}
